\section{Preliminaries and Problem Formulation}
\label{sec:prelim}

\subsection{Attribution Methods as Maps of the Parameters}
We treat every post-hoc attribution as a deterministic map from an input
$\bx$ and the parameters $\btheta$ to a heatmap $\bvphi(\bx;\btheta)\in\R^{d}$.
For the target class $c=\argmax_k S_k(\bx;\btheta)$ we consider four
representative methods spanning the gradient/perturbation and
local/class-discriminative axes.

\paragraph{Vanilla saliency \cite{simonyan2014saliency}}
\begin{equation}
\bvphi^{\mathrm{S}}(\bx;\btheta)=\bigl|\nabla_{\bx}S_c(\bx;\btheta)\bigr|,
\label{eq:sal}
\end{equation}
the element-wise magnitude of the input gradient.

\paragraph{Integrated Gradients \cite{sundararajan2017ig}}
With baseline $\bx'$ (here $\bx'=\bm 0$),
\begin{equation}
\varphi^{\mathrm{IG}}_i(\bx;\btheta)=(x_i-x'_i)\!\int_{0}^{1}\!
\frac{\partial S_c\!\bigl(\bx'+\alpha(\bx-\bx');\btheta\bigr)}{\partial x_i}\,
\mathrm d\alpha,
\label{eq:ig}
\end{equation}
approximated by an $n$-point Riemann sum. IG satisfies \emph{completeness},
$\sum_i\varphi^{\mathrm{IG}}_i=S_c(\bx)-S_c(\bx')$.

\paragraph{Grad-CAM \cite{selvaraju2017gradcam}}
Let $\bm A^{k}\in\R^{u\times v}$ be the $k$-th feature map of the last
convolutional block and
$\alpha_c^{k}=\tfrac{1}{uv}\sum_{i,j}\partial S_c/\partial A^{k}_{ij}$. Then
\begin{equation}
\bvphi^{\mathrm{GC}}(\bx;\btheta)=\ReLU\!\Bigl(\textstyle\sum_{k}\alpha_c^{k}
\bm A^{k}\Bigr),
\label{eq:gradcam}
\end{equation}
bilinearly up-sampled to $H\times W$.

\paragraph{Occlusion \cite{zeiler2014visualizing}}
For a mask $\bm m_p$ over region $p$,
\begin{equation}
\varphi^{\mathrm{Occ}}_p(\bx;\btheta)=\bigl[S_c(\bx;\btheta)-
S_c(\bx\odot(1-\bm m_p);\btheta)\bigr]_{+}.
\label{eq:occ}
\end{equation}

All maps are min--max normalised to $[0,1]$ so that methods of differing native
scale are compared on common footing.

\subsection{Compression Operators}
\begin{definition}[Uniform symmetric quantization]
For a weight tensor $\bW$ with dynamic range $R=\max|\bW|$ and bit-width $b$,
let $q_{\max}=2^{\,b-1}-1$ and step $\Delta_q=R/q_{\max}$. Then
\begin{equation}
Q_b(\bW)=\Delta_q\cdot\mathrm{clip}\!\Bigl(\Bigl\lfloor\tfrac{\bW}{\Delta_q}
\Bigr\rceil,-q_{\max}-1,q_{\max}\Bigr),
\label{eq:quant}
\end{equation}
where $\lfloor\cdot\rceil$ is round-to-nearest.
\end{definition}

\begin{definition}[Global magnitude pruning]
Given target sparsity $s\in[0,1)$, let $\tau_s$ be the $s$-quantile of
$\{|\theta_i|\}$. Then $P_s(\btheta)=\btheta\odot\bm 1[\,|\btheta|>\tau_s\,]$.
\end{definition}

Both operators yield $\tilde\btheta=\Cop(\btheta)$ with
$\Delta\btheta=\tilde\btheta-\btheta$; quantization perturbs \emph{every}
weight by a bounded amount, whereas pruning perturbs a \emph{subset} by their
full value.

\subsection{Fidelity, Localisation and Faithfulness}
We quantify the effect of compression on explanations through three families
of metrics. Let $\bvphi=\bvphi(\bx;\btheta)$ and $\tilde\bvphi=\bvphi(\bx;
\tilde\btheta)$.

\emph{(i) Model-vs-model fidelity} measures how much the explanation
\emph{changed}:
\begin{align}
F_{\rho}&=\corr_{\mathrm{Spearman}}(\bvphi,\tilde\bvphi),\quad
F_{r}=\corr_{\mathrm{Pearson}}(\bvphi,\tilde\bvphi),\\
F_{\mathrm{IoU}}&=\frac{|\mathcal T_k(\bvphi)\cap\mathcal T_k(\tilde\bvphi)|}
{|\mathcal T_k(\bvphi)\cup\mathcal T_k(\tilde\bvphi)|},\quad
F_{\mathrm{SSIM}}=\mathrm{SSIM}(\bvphi,\tilde\bvphi),
\end{align}
where $\mathcal T_k(\cdot)$ is the set of top-$k$ salient pixels
($k=20\%$). $F_\rho=1$ means the ranking of pixel importances is unchanged.

\emph{(ii) Localisation vs.\ ground truth} exploits the fact that our benchmark
provides a binary object mask $\bm M$: for the compressed map,
\begin{equation}
\mathrm{IoU}_{\mathrm{GT}}=\frac{|\mathcal T_k(\tilde\bvphi)\cap\bm M|}
{|\mathcal T_k(\tilde\bvphi)\cup\bm M|},\quad
\mathrm{PG}=\bm 1\!\bigl[\argmax\tilde\bvphi\in\bm M\bigr],
\end{equation}
the top-$k$ overlap with, and pointing-game hit rate inside, the true object.

\emph{(iii) Faithfulness} measures whether the map identifies pixels the model
actually uses, via the deletion curve \cite{petsiuk2018rise}: pixels are removed
in descending saliency order and the area under the target-probability curve is
recorded; a \emph{lower} deletion AUC indicates a more faithful map. Two
complementary faithfulness summaries, standard in sequence and regression
explanation studies \cite{samek2017evaluating}, are also supported by our
protocol: the \emph{Area Over the Perturbation Curve} (AOPC), the mean drop in
target probability as the most-relevant regions are progressively removed
(higher is more faithful), and the \emph{Ablation Percentage Threshold} (APT),
the fraction of features that must be removed to drive the prediction below a
fixed confidence.

\emph{(iv) Sample-level consistency.} Metrics (i)--(iii) report dataset averages,
which can mask heavy-tailed behaviour across individual inputs. To characterise
the \emph{distribution} of explanation quality we additionally report the
skewness ($\mathrm{AUC\ Skew}$) and excess kurtosis ($\mathrm{AUC\ Kurt}$) of the
per-sample score-drop AUC. A near-zero skew and kurtosis indicate that
compression degrades explanations uniformly across samples, whereas large values
flag a subpopulation of inputs whose explanations collapse well before the
dataset mean would suggest.

\begin{definition}[Explanation fidelity of a compression]
For a fixed method and metric $F\!\in\![0,1]$ and a distribution $\mathcal D$
of inputs, the fidelity of compression level $\ell$ is
\begin{equation}
\mathcal F(\ell)=\E_{\bx\sim\mathcal D}\,F\!\bigl(\bvphi(\bx;\btheta),
\bvphi(\bx;\Cop_\ell(\btheta))\bigr).
\end{equation}
\end{definition}

The central objects of study are the profiles $\mathcal F(b)$ (quantization)
and $\mathcal F(s)$ (pruning), and the design of a compression operator that
keeps $\mathcal F$ high at a given budget.
